\section{Modellierung mechatronischer Systeme}

\subsection{Systematik der Modellbildung}
\begin{tcolorbox}
\boxtitle{Vier Phasen der Modellierung}
\begin{enumerate}
  \item \textbf{Ziel bestimmen:} Reglerentwurf, Simulation, Belastungsanalyse? Statisch oder dynamisch? Geforderte Modellgenauigkeit?
  \item \textbf{Modellstruktur festlegen:} Ein- und Ausgangsgrößen definieren, Systemgrenzen ziehen, Zerlegung in Teilsysteme (elektrisch, mechanisch, thermisch).
  \item \textbf{Parameter bestimmen:} Aus Datenblättern, Messungen im Arbeitspunkt oder Identifikationsverfahren.
  \item \textbf{Modellvalidierung:} Plausibilitätsprüfungen (Einheiten, Vorzeichen) und Vergleich der Simulation mit realen Messdaten.
\end{enumerate}
\end{tcolorbox}

\subsection{Modellierung am Beispiel Drosselklappe}

\begin{tcolorbox}
\boxtitle{1. Elektrisches Teilsystem (Gleichstrommotor)}
Maschengleichung des Ankerkreises:
\[
u_A(t) = R_A \cdot i_A(t) + L_A \cdot \frac{di_A(t)}{dt} + u_{\text{EMK}}(t)
\]
Elektromotorische Gegenkraft (EMK):
\[
u_{\text{EMK}}(t) = k_{\text{EMK}} \cdot \omega_M(t) = k_{\text{EMK}} \cdot k_G \cdot \dot{\varphi}(t)
\]
Elektromagnetisches Motormoment:
\[
M_{\text{Mot}}(t) = k_M \cdot i_A(t)
\]
\end{tcolorbox}

\begin{tcolorbox}
\boxtitle{2. Mechanisches Teilsystem (Klappenmechanik)}
Drallsatz um die Drehachse der Drosselklappe:
\[
J \cdot \ddot{\varphi}(t) + d \cdot \dot{\varphi}(t) + M_F(\varphi) + M_{\text{Reib}}(\dot{\varphi}) = M_{\text{Klappe}}(t)
\]
mit dem Antriebsmoment $M_{\text{Klappe}} = k_G \cdot M_{\text{Mot}}$:
\begin{itemize}
  \item $J$: Gesamt-Massenträgheitsmoment auf Klappenachse bezogen.
  \item $d$: Viskoser Dämpfungsbeiwert.
  \item $M_F(\varphi) = c_F \cdot (\varphi(t) - \varphi_0)$: Rückstellmoment der Rückholfeder (Fail-Safe Ruhelage $\varphi_0$).
  \item $M_{\text{Reib}}$: Reibung (Coulomb-Gleitreibung und Haftreibung).
  \item $k_G$: Übersetzungsverhältnis des Getriebes ($k_G = \omega_M / \dot{\varphi}$).
\end{itemize}
\end{tcolorbox}

\begin{tcolorbox}
\boxtitle{3. Modellvereinfachungen für den Reglerentwurf}
\begin{itemize}
  \item \textbf{Elektrische Zeitkonstante:} Da $T_{\text{el}} = \frac{L_A}{R_A} \ll T_{\text{mech}}$, wird die Ankerinduktivität für den linearen Reglerentwurf meist vernachlässigt ($L_A \approx 0$):
  \[
  i_A(t) \approx \frac{u_A(t) - k_{\text{EMK}} \cdot k_G \cdot \dot{\varphi}(t)}{R_A}
  \]
  \item \textbf{Vernachlässigung von Feder und Reibung:}
  Für schnelle Lageänderungen dominiert das Trägheitsmoment $J \ddot{\varphi}$. Das System nähert sich einem idealen $IT_1$-Verhalten von Eingangsspannung $u_A$ zu Klappenwinkel $\varphi$:
  \[
  F_S(s) = \frac{\Phi(s)}{U_A(s)} \approx \frac{K}{s (1 + s T_{\text{mech}})}
  \]
  mit $T_{\text{mech}} = \frac{J \cdot R_A}{d \cdot R_A + k_M \cdot k_{\text{EMK}} \cdot k_G^2}$.
\end{itemize}
\end{tcolorbox}

\begin{tcolorbox}
\boxtitle{Plausibilitätsprüfungen}
\begin{itemize}
  \item Vorzeichen der Kennlinien prüfen (z.\,B. Reibmoment wirkt immer der Bewegungsrichtung entgegen: $\text{sign}(\dot{\varphi})$).
  \item Federmoment muss im stromlosen Zustand größer sein als die Reibung ($M_F > M_{\text{Reib}}$), damit Fail-Safe gesichert ist.
  \item Maximales Motormoment muss Feder- und Reibmoment sicher überwinden: $k_M \cdot I_{\max} \cdot k_G > M_F + M_{\text{Reib}}$.
\end{itemize}
\end{tcolorbox}
